% Part: proof-theory
% Chapter: proof-search
% Section: search-algorithm

\documentclass[../../../include/open-logic-section]{subfiles}

\begin{document}

\olfileid{pt}{ps}{sal}

\olsection{\Log{G3c} માટે સાબિતી-શોધ અલ્ગોરિધમ}

\Log{G3c} માટે
સાબિતી-શોધ
અલ્ગોરિધમ વર્ણવીએ.
શક્ય એવો
આ એકમાત્ર
અલ્ગોરિધમ નથી,
અને સૌથી
કાર્યક્ષમ પણ
નથી. પણ તે
સાચો છે:
સિક્વન્ટ
$\Gamma \Sequent \Delta$ની
!!a{proof}
હોય, તો
તે આખરે
થોભીને
!!a{proof}
શોધશે. તે
અચૂક પણ
છે: !!a{proof}
મળ્યા વિના
થોભે અથવા
કદી થોભે
જ નહીં,
તો
$\Gamma
\Sequent \Delta$
શરૂઆતથી જ
અપ્રામાણ્ય
છે.

અલ્ગોરિધમનું
આવશ્યક લક્ષણ
એ છે કે
$\Gamma \Sequent \Delta$
અથવા શોધ
દરમિયાન મળતા
સિક્વન્ટનું
દરેક
!!{formula}
જરૂર પડે
તેટલી વાર
``ઘટાડાય'':
એટલે તેને
ઊલટી દિશામાં
લાગુ કરેલા
અનુમાન-નિયમનું
મુખ્ય સૂત્ર
ગણાય. આ
લક્ષણને
\emph{ન્યાયીપણું}
કહીએ છીએ.

અલ્ગોરિધમ
તબક્કાવાર
કામ કરે છે.
દરેક તબક્કાનું
પરિણામ
સિક્વન્ટોનું
વૃક્ષ છે;
આગલા તબક્કામાં
પહેલેથી સ્વયંસિદ્ધ
ન હોય એવા
દરેક સૌથી
ઉપરના સિક્વન્ટમાં
!!a{formula}
ઘટાડીને
નવું વૃક્ષ
બનાવાય છે.

તબક્કા~$0$માં
$\Gamma \Sequent \Delta$થી
શરૂ કરીએ.
ન્યાયીપણું
સુનિશ્ચિત કરવા
$\Gamma$
અને~$\Delta$નાં
!!{formula}sને
અનુક્રમણિકાંક
તરીકે સંખ્યાઓ
આપીએ, જેથી
અલગ
!!{formula}sને
અલગ સંખ્યા
મળે. દાખલા
તરીકે ડાબેથી
જમણે
!!{formula}sની
ગણતરી કરો.
સંખ્યાઓ
ઉપરિચિહ્ન
તરીકે
લખીએ.
$!A, !B \Sequent !C$ની
!!a{proof}
શોધતા હોઈએ,
તો તબક્કા~$0$નું
પરિણામ
$!A^1, !B^2 \Sequent !C^3$
છે. આવા
અનુક્રમણિકાંકવાળા
સિક્વન્ટમાં
ઉપરિચિહ્ન
તરીકે આવતો
સૌથી મોટો~$n$
તેનો
\emph{મહત્તમ
અનુક્રમણિકાંક}
છે (દાખલામાં~$3$).

સિક્વન્ટમાં
પરિમાણકાર
હોય, તો જમણે
$\lexists[x][!A(x)]$
અથવા ડાબે
$\lforall[x][!A(x)]$
ઘટાડતાં કયાં
પદો વાપરવાં
તે પણ જાણવું
પડે. માટે
ફક્ત
!!{constant}s
$\Obj c_0$,
$\Obj c_1$,
$\Obj c_2$,
\dots, અને
$\Gamma$
તથા~$\Delta$માં
આવતાં ફલન-ચિહ્નો
વડે બનેલાં
પદો જરૂરી
છે. આ
બધાં પદો
કોઈ સ્થિર
ક્રમમાં ગોઠવેલાં
છે એમ
માનીએ, જેથી
દાખલા તરીકે
``$\Gamma \Sequent \Delta$માં
ન આવતો
પ્રથમ
!!{constant}''
કહી શકાય.
અલ્ગોરિધમના
વૃક્ષમાં દરેક
અનુક્રમણિકાંકવાળા
સિક્વન્ટ સાથે
!!{constant}sનો
ગણ સંકળાયેલો
છે. તબક્કા~$0$ના
સિક્વન્ટને
તેમાં આવતા
બધા
!!{constant}sનો
ગણ આપીએ
(એવા !!{constant}s
હોય તો);
એકેય !!{constant}
ન
હોય તો
$\{\Obj c_0\}$
આપીએ.

માનો કે
તબક્કા~$n$એ
અલ્ગોરિધમને
અનુક્રમણિકાંકવાળા
સિક્વન્ટોનું
વૃક્ષ અને
તેમની સાથે
!!{constant}sના
ગણો મળ્યા
છે. તબક્કા~$n+1$એ
દરેક સૌથી
ઉપરના સિક્વન્ટ
માટે આ
કરીએ:

!!a{formula}~$!A$
સિક્વન્ટની
ડાબી અને
જમણી બંને
બાજુ હોય,
અથવા ડાબે
$\lfalse$
હોય, તો
કંઈ કરતા
નથી: આવા
સિક્વન્ટોની
\Log{G3c}માં
!!{proof}s
છે, એટલે
આ સૌથી
ઉપરના સિક્વન્ટ
માટે કામ
પૂરું થયું.

સિક્વન્ટમાં
બિનઆણ્વિક
!!{formula}
ન હોય,
તો અલ્ગોરિધમ
અટકાવો.
$\Gamma \Sequent \Delta$ની
સાબિતી નથી.

નહિતર સૌથી
નાના
અનુક્રમણિકાંક~$i$
વાળું
બિનઆણ્વિક
!!{formula}~$!A$
પસંદ કરો.
ત્યારે સંબંધિત
સૌથી ઉપરનો
સિક્વન્ટ
$!A^i, \Pi
\Sequent \Lambda$
અથવા
$\Pi \Sequent \Lambda, !A^i$
છે.
$k$ એ
આ સિક્વન્ટના
મહત્તમ
અનુક્રમણિકાંક
કરતાં એક
વધુ, એટલે
આગલો ન
વપરાયેલો
અનુક્રમણિકાંક,
અને~$C$
તેને સોંપેલો
!!{constant}sનો
ગણ હોય.
\sourcecorrection{OLMD-275}{મૂળ $k$ને હાલનો મહત્તમ અનુક્રમણિકાંક કહે છે, પરંતુ નીચેના નિયમો નવા સૂત્રોને $k$ અને $k+1$ આપે છે. તેમની તાજગી તથા ન્યાયી પસંદગી માટે $k$ હાલના મહત્તમથી એક વધુ હોવો જોઈએ.}

$!A^i$ને
``ઘટાડીએ'':
એટલે~$!A$ના
!!{main operator}
અને~$!A^i$
ડાબે છે
કે જમણે
તે મુજબ
નક્કી થતા
અનુમાન-નિયમથી
નવા સૌથી
ઉપરના સિક્વન્ટો
બનાવીએ.

\begin{enumerate}
  \item $!A \ident !B \land !C$ હોય અને ડાબે આવતું હોય:
  $(!B \land !C)^i, \Pi \Sequent \Lambda$ની ઉપર એક નવો
  સૌથી ઉપરનો સિક્વન્ટ ઉમેરો:
  \[
  \Axiom$!B^k, !C^{k+1}, \Pi \fCenter \Lambda$
  \RightLabel{\LeftR{\land}}
  \UnaryInf$(!B \land !C)^i, \Pi \fCenter \Lambda$
  \DisplayProof
  \]
  નવા સિક્વન્ટને સોંપેલો !!{constant}sનો ગણ~$C$ છે.
  \item $!A \ident !B \land !C$ હોય અને જમણે આવતું હોય:
  $\Pi \Sequent \Lambda, (!B \land !C)^i$ની ઉપર બે નવા
  સૌથી ઉપરના સિક્વન્ટો ઉમેરો:
  \[
  \Axiom$\Pi \fCenter \Lambda, !B^k$
  \Axiom$\Pi \fCenter \Lambda, !C^k$
  \RightLabel{\RightR{\land}}
  \BinaryInf$\Pi \fCenter \Lambda, (!B \land !C)^i$
  \DisplayProof
  \]
  નવા સિક્વન્ટોને સોંપેલો !!{constant}sનો ગણ~$C$ છે.

  \item $!A \ident \lnot !B$, $!A \ident !B \lor !C$,
  અને $!A \ident !B \lif !C$ના કિસ્સા સમાન છે અને
  અભ્યાસ માટે છોડ્યા છે.

  \item $!A \ident \lexists[x][!B(x)]$ હોય અને ડાબે આવતું હોય:
  $\lexists[x][!B(x)]^i, \Pi \Sequent \Lambda$ની ઉપર એક નવો
  સૌથી ઉપરનો સિક્વન્ટ ઉમેરો:
  \[
  \Axiom$!B(c)^k, \Pi \fCenter \Lambda$
  \RightLabel{\LeftR{\lexists}}
  \UnaryInf$\lexists[x][!B(x)]^i, \Pi \fCenter \Lambda$
  \DisplayProof
  \]
  અહીં $c$ એ~$C$માં ન હોય એવો પ્રથમ !!{constant} છે.
  નવા સિક્વન્ટને સોંપેલો !!{constant}sનો ગણ $C \cup \{c\}$ છે.
  \item $!A \ident \lexists[x][!B(x)]$ હોય અને જમણે આવતું હોય:
  $\Pi \Sequent \Lambda, \lexists[x][!B(x)]^i$ની ઉપર એક નવો
  સૌથી ઉપરનો સિક્વન્ટ ઉમેરો:
  \sourcecorrection{OLMD-277}{મૂળ ગદ્યમાં આ નિષ્કર્ષના અસ્તિત્વ-સૂત્રનો અનુક્રમણિકાંક છૂટી ગયો છે; નીચેના નિયમવૃક્ષ મુજબ અહીં $i$ ઉમેર્યો છે.}
  \[
  \Axiom$\Pi \fCenter \Lambda, \lexists[x][!B(x)]^k, !B(t)^{k+1}$
  \RightLabel{\RightR{\lexists}}
  \UnaryInf$\Pi \fCenter \Lambda, \lexists[x][!B(x)]^i$
  \DisplayProof
  \]
  અહીં $t$ એ ફક્ત~$C$માંના !!{constant}s ધરાવતું એવું
  પ્રથમ પદ છે જે આ સિક્વન્ટની નીચે જમણી બાજુના
  $\lexists[x][!B(x)]$ના ઘટાડામાં અગાઉ વપરાયું નથી
  (અથવા આવાં બધાં પદો વપરાઈ ગયાં હોય તો $\Obj c_0$).
  નવા સિક્વન્ટને સોંપેલો !!{constant}sનો ગણ $C$ છે.
  \sourcecorrection{OLMD-276}{મૂળ નિયમવૃક્ષમાં જમણી બાજુના અસ્તિત્વ-સૂત્ર માટે ભૂલથી ડાબો નિયમ \LeftR{\lexists} લખાયો છે; તેને યોગ્ય \RightR{\lexists} કર્યો છે.}
  \item $!A \ident \lforall[x][!B(x)]$ના કિસ્સા સમાન છે અને
  અભ્યાસ માટે છોડ્યા છે.
\end{enumerate}

તબક્કા~$n+1$એ કોઈ પણ સૌથી ઉપરના સિક્વન્ટમાં નવો
સિક્વન્ટ ઉમેરાયો ન હોય, તો અલ્ગોરિધમ સફળતાપૂર્વક
થોભે છે. આ કિસ્સામાં તબક્કા~$n+1$ના વૃક્ષમાંથી
બધા અનુક્રમણિકાંક ભૂંસીને $\Gamma \Sequent \Delta$ની
!!a{proof} મળે છે. સૌથી ઉપરના બધા સિક્વન્ટો
$!A, \Pi \Sequent \Lambda, !A$ અથવા
$\lfalse, \Pi \Sequent \Lambda$ સ્વરૂપના છે.
બધાં અનુમાનો \Log{G3c}નાં યોગ્ય અનુમાનો છે.
વિધાનતર્કીય અનુમાનો માટે આ સ્પષ્ટ છે.
નોંધો કે દરેક પગલે સિક્વન્ટને સોંપેલા
!!{constant}sના ગણ~$C$માં તે સિક્વન્ટમાં આવતા
બધા !!{constant}s હોય છે. તેથી
\LeftR{\lexists} અને \RightR{\lforall} વાપરીને
કરેલા ઘટાડામાં સ્વનિર્ધારક ચલની શરત પૂરી થાય છે:
સ્વનિર્ધારક ચલ~$c$ નિષ્કર્ષને સોંપેલા ગણ~$C$માં
ન હોય એવો !!a{constant} હતો, એટલે નિષ્કર્ષમાં
એટલે $c$ નિષ્કર્ષમાં
પણ આવતો નહોતો. આમ:

\begin{prop}
\Log{G3c} માટેનો સાબિતી-શોધ અલ્ગોરિધમ સાચો છે:
તે સફળતાપૂર્વક થોભે, તો શરૂઆતના
સિક્વન્ટ~$\Gamma \Sequent \Delta$ની !!a{proof} છે.
\end{prop}





\end{document}
